\section{Introduction}\label{sec:introduction}

The question (T) asks whether one fixed deterministic multitape Turing machine
can transpose all square matrices in time linear in their row-major encoding
length. The input is read-only; work tapes are read/write; output is append-only
and cannot be read. A mathematical decoder may inspect records in any order,
whereas the machine must obtain its operands with the actual tape heads.

The seventeenth paper identifies a gap between restrictions on decoding
rules and restrictions on complete tape computations
\cite[Sections~3 and~7.4]{SeventeenthPaper2026}. Its order and symmetry
theorems apply to structured one-hop decoders. Obtaining such structure,
with its advice and intermediate states paid for, from arbitrary native
executions remains open. We work on the execution side of this gap:
which records must carry information, where must that information pass,
and how much can fixed local descriptions explain?

\paragraph{Original adaptive executions.}
Earlier macroscopic-crossing results constrain oblivious executions and
sufficiently heavy original traces. Their general adaptive application uses
a low-trace simulator whose dependency graph differs from the original
machine's graph \cite{MacroscopicCrossings2026}. Our direct replay instead
pays for selected intervals of the original execution.

Fix the number of work tapes, the alphabet and a linear time constant.
There are positive constants $\varepsilon,\gamma,\eta$ such that fewer than
$2^{(1-\eta)P}$ length-$P$ inputs are both timely and correctly transposed
and admit a free work head leaving at most $\gamma P$ returns of excursion
greater than $\varepsilon P$ on the other heads. The threshold is uniform
over all positive entry widths. Consequently two heads make linearly many
macroscopic returns on every other timely correct input. The exception
count is relative to the ambient cube; it is not a conditional proportion
of an arbitrarily small success set.

The proof retains every visit to a spatial slab of the free head. Its
memory persists through gaps between retained runs. On another tape, a
first read outside the paid entry window either reads a virgin cell or has
an actual long-return predecessor. One joint sparse event/value record
supplies these symbols. The replayed correct outputs recover raw bits
omitted from the message.

A second replay retains causal fragments rather than every step in a slab.
It pays complete cut values and processes intervening writes chronologically.
This gives simultaneous expansion for arbitrary macroscopic raw demand sets
and, after localization, expansion to off-carrier long-return predecessors.
The localized result concerns the output sets supplied by the slab geometry.

Two further conclusions use the existing aligned simulator of
\cite{HighRateStorage2026}. We prove an actual path and cone correspondence,
paying guards for the changed clock, head positions and delayed outputs.
It transfers expansion from \emph{every} large output set to original
crossings and gives storage-capacity certificates on the original entry
relation. Their asymptotic statements use canonical widths. The finite
eroded-input cut invoice and the direct replay results hold at every
positive width. No common original trace is assumed in either route.

\paragraph{What stored data must do.}
A one-work-tape simulation handles runs in which at every configuration
at most one work tape has more than $m$ nonblank cells. The carrier may
change arbitrarily often: sparse sketches and a lazily reset visited
interval preserve the computation through a handoff. Combining the
simulation with the one-head bound forces two simultaneous large stores.
A linear-overhead redundant-archive wrapper shows why this conclusion
does not imply independence of their contents.

The extensions then test more informative descriptions of those stores.
Fixed getters reconstruct incoming values from continuous raw windows,
paid observer state, one shared label and the strictly prior correct-output
string. Shared spatial checkpoints supply such getters for deterministic
scans. Changing-carrier and historical-cache arguments compare whole words
with fixed locally readable reference families and force long secondary
modification episodes. Persistence estimates quantify marginal uncertainty,
including uncertainty after conditioning on the emitted prefix.
The historical-feedback appendix measures dated work-block reuse through
actual halt. A cheap copying control bounds its interpretation: a large
historical statistic alone does not account for the cost of transposition.

\paragraph{Proofs and formal status.}
Each numbered result identifies its formal status. The replay, original
return and cut bounds, coherent-cluster extension, getters, paid loader,
complete block-program lemma and actual simulator transport are proved
in Lean. The carrier, catalogue, eraser, archive and physical monitor
constructions have complete written proofs but retain native formalization
obligations. Their checked consequences explicitly assume the missing
construction bounds. No additional axiom supplies them.
The Lean declarations are authoritative; the English outlines are reading
aids. \Cref{sec:formalization} states the precise boundary.

\paragraph{Relation to earlier work and reading order.}
The native locality mechanism comes from \cite{FirstPaper2026}, and the
finite arbitrary-width slab geometry from \cite{MacroscopicCrossings2026}.
The additional work is the paid replay across different original executions,
the actual simulator pullback, and the carrier and reference constructions.
We reuse standard counting and entropy inequalities. No external-priority
claim is made for the combined statements.

For the direct native bounds, read \cref{sec:replay,sec:returns,sec:fragments}.
For the every-output crossing and entry-capacity route, read
\cref{sec:accounting,sec:transfer}. \Cref{sec:storage,sec:limitations}
give the simultaneous-store condition and its redundancy control.
The other appendices develop coherent clusters, catalogues, local getters,
reference changes, persistence and feedback. \Cref{sec:open} specifies
which additional statements would turn these requirements into an
impossibility proof.
